\documentclass[12pt]{cwru_thesis}
\usepackage[hyphens]{url}
\usepackage{lipsum}
\usepackage{graphicx}
\usepackage{hyperref}
\hypersetup{colorlinks=true, linkcolor=black, filecolor=black, urlcolor=black, citecolor=black}
\urlstyle{same}
\usepackage[acronym,toc]{glossaries}
\usepackage[intoc]{nomencl}
\renewcommand{\nomname}{List of Symbols}

\usepackage{mathptmx}
% \usepackage[T1]{fontenc}
% \usepackage[utf8]{inputenc}
\usepackage{amsfonts, amsmath, amsthm, amssymb}
\usepackage{newtxtext,newtxmath}

% my packages
\usepackage{minted}
% \usepackage{unicode-math}
% \usepackage{fontspec}
% \usepackage{polyglossia}

\newcommand{\dif}{\mathrm{d}}
\newcommand{\Dif}{\mathrm{D}}
\renewcommand{\vec}[1]{\ensuremath{\underline{#1}}}
\newcommand{\grad}{\underline{\nabla}}
\newcommand{\curl}{\underline{\nabla} \times}
\newcommand{\lap}{\underline{\nabla}^2}
\renewcommand{\div}{\underline{\nabla} \cdot}
\newcommand{\paren}[1]{\left( {#1} \right)}
\newcommand{\norm}[1]{\lVert#1\rVert}

% my own commands
\newcommand{\code}[1]{
    \inputminted[
        mathescape=true,
        linenos=true,
        numbersep=3pt,
        frame=single,
        numberblanklines=true,
        breaklines=true,
        breaksymbol=,
        framesep=2mm,
        baselinestretch=1
        ]{julia}{#1}
    }

% Must use biblatex to produce the Published Contents and Contributions, per-chapter bibliography (if desired), etc.
\usepackage[
    backend=biber,natbib,style=numeric-comp,maxbibnames=99
]{biblatex}

% Name of your .bib file(s)
\addbibresource{example.bib}

\makeglossaries
%%%%%%%%%%%%%%%%%%%%%%%%%%%%%%%%%%%%%%%%%%%%%%%%%%%%%%%%%%%%%%%%%%%%%%%%%%%%%%%%%%%%%%%
%Glossary entries
\newglossaryentry{galaxy}
{
    name=galaxy,
    description={A system of stars independent from all other systems}
}
\newglossaryentry{asteroid}
{
    name=asteroid,
    description={A very small planet ranging from 1,000 km to less than one km in diameter. Asteroids are found commonly around other larger planets}
}

%Acronyms to include in the list of acronyms
\newacronym{leo}{LEO}{Low Earth Orbit}
%%%%%%%%%%%%%%%%%%%%%%%%%%%%%%%%%%%%%%%%%%%%%%%%%%%%%%%%%%%%%%%%%%%%%%%%%%%%%%%%%%%%%%%

\makenomenclature
%%%%%%%%%%%%%%%%%%%%%%%%%%%%%%%%%%%%%%%%%%%%%%%%%%%%%%%%%%%%%%%%%%%%%%%%%%%%%%%%%%%%%%%
%Nomenclature entries
\nomenclature{\(\alpha\)}{Angle of attack (degrees)}
\nomenclature{\(c\)}{Speed of light in a vacuum (m/s)}
\nomenclature{\(\mathbb{R}\)}{Real numbers}
%%%%%%%%%%%%%%%%%%%%%%%%%%%%%%%%%%%%%%%%%%%%%%%%%%%%%%%%%%%%%%%%%%%%%%%%%%%%%%%%%%%%%%%

\begin{document}
% \pagenumbering{roman}

\title{Study of dynamical systems using the Julia programming language} %Title of your thesis in title case
\author{Jonáš Koziorek} %Your name

\degreeaward{Bachelor of mathematics} % Degree to be awarded
\department{Department of Applied Mathematics}
\university{VSB - Technical University of Ostrava}    % Institution name  
\unilogo{vsb_logo.eps}                                 % Institution logo
\defendmonth{May, 2024}      % Graduation month and year
\defenddate{30.4.2024}     % Date of thesis defense

%Committee Member names. If you have a different number of committee members for your defense, you will need to edit lines 183-190 of cwru_thesis.cls accordingly.
\committeeChair{Marek Lampart} %Committee Chair's name
\committeeOne{Jiří Tomčala} %Committee member #1's name
\committeeTwo{[Committee Member]} %Committee member #2's name
\committeeThree{[Committee Member]} %Committee member #3's name

\maketitle[logo]

\begin{dedication} 	 
    This thesis is dedicated to ..... to be added.
\end{dedication}

\begin{KeepFromToc}
  \tableofcontents
\end{KeepFromToc}
\listoftables
\listoffigures

\begin{preface}
    preface ... to be added
\end{preface}

\begin{acknowledgements}
    I would like to thank ... to be added.
\end{acknowledgements}

%If you have acroynms you wish to define, include this line. If not, you may delete it. See https://www.overleaf.com/learn/latex/Glossaries#Acronyms for more information about how to use Acronyms.
\printglossary[type=\acronymtype, title=List of Abbreviations]

%If you have a nomenclature section for defining symbols, include this line. If not, you may delete it. See https://www.overleaf.com/learn/latex/Nomenclatures for more information about how to use Nomenclatures
\printnomenclature

\begin{abstract}
    abstract to be added
\end{abstract}

\mainmatter

\setcounter{secnumdepth}{2}

\chapter{Introduction} \label{chap:intro}
Start off all chapters with \verb|chapter|. \verb|\extrachapter| will give you an unnumbered chapter that's added to the Table of Contents. Labeling is also good practice. That way you can reference sections and chapters from anywhere. For instance, I can reference Sec.~\ref{sec:exsection} from here.

Here's an example of a citation \citep{GMP81}. Here's another \citep{PP98}. These will appear in the big bibliography at the end of the thesis.

If you're new to \LaTeX{} and would like to begin by learning the basics, please see our free online course available at:\\ \url{https://www.overleaf.com/latex/learn/free-online-introduction-to-latex-part-1} \index{LaTeX@\LaTeX}

You can define terms and/or symbols in a glossary \index{glossary} as you talk about key terms in your thesis. So what's a \gls{galaxy}? These can also include acronyms, such as \acrlong{leo}, which is abbreviated as \acrshort{leo}. You should always define acronyms at the point of first use in the text, in addition to including them in the glossary.

\code{codes/havok.jl}

\section{This is a Section} \label{sec:exsection}
\lipsum[1-2]

Check this out! I can reference a figure number in the text of the document by referring to Fig.~\ref{fig:logo}.

\begin{figure}[hbt!]
    \centering
    \includegraphics[width=.5\textwidth]{vsb_logo.eps}
    \caption{This is a figure. \label{fig:logo}}
\end{figure}

\subsection{This is a Subsection} \label{subsec:exsubsec}

\begin{table}[hbt!]
\centering
\begin{tabular}{ll}
\hline
Area & Count\\
\hline
North & 100\\
South & 200\\
East & 80\\
West & 140\\
\hline
\end{tabular}
\caption{This is a table. \label{tab:sample}}
\end{table}

\lipsum[3] 

I can add another glossary entry wherever I like, and I can make it capitalized by using the macro command with a capital letter like this: \Gls{asteroid}. The quadratic formula is given in Eq.~\eqref{eq:quadratic}.
\begin{equation}
    x=\frac{-B \pm \sqrt{B^2 - 4AC}}{2A}
    \label{eq:quadratic}
\end{equation}

\lipsum[4-5] 

\section{This is Another Section}
\lipsum[6-7] 

\chapter{This is the Second Chapter}
\lipsum[8-9]


\chapter{This is the Third Chapter}
See? I can reference the first chapter from all the way down here. I think the first chapter was Chapter~\ref{chap:intro}.
\chapter{This is the Fourth Chapter}
\chapter{This is the Fifth Chapter}
\chapter{This is the Sixth Chapter}
\chapter{This is the Seventh Chapter}
\chapter{This is the Eighth Chapter}


%If you have appendices to your thesis, place them here
\appendix

\chapter{Questionnaire}

\printbibliography[heading=bibintoc]

\end{document}
